% --- Archon physics formalization source begin ---
% archon:physics
% archon:covers PhyXMiniProblems/problem_phyx_mini_0178.lean
% archon:source-report reports/phyx_mini/problem_phyx_mini_0178.source.json

\chapter{Physics problem phyx\_mini\_0178}
\label{ch:PhyXMiniProblems_problem_phyx_mini_0178}

\paragraph{Problem source.}
\#\# Physical scenario

Figure shows a standing wave that is oscillating at frequency \$f\_0\$.

\#\# Question

If the tension in the string is increased by a factor of four, for what frequency, in terms of \$f\_0\$, will the string continue to oscillate as a standing wave with four antinodes?

\#\# Answer choices

- A: A: \$\textbackslash{}sqrt\{2\} f\_0\$

- B: B: \$f\_0\$

- C: C: \$\textbackslash{}frac\{f\_0\}\{2\}\$

- D: D: \$2 f\_0\$

\#\# Figure caption (auxiliary; use the image as primary evidence)

The image depicts two parallel sinusoidal waves positioned horizontally across the center. The waves are blue, and their crests and troughs are symmetrically aligned, creating a mirrored effect. On both sides, there are vertical gradient bars with a light-to-dark pattern, giving an impression of fading or shadowing. The waves seem to start and end at the points where these vertical bars are located. There is no text present in the image.

\#\# Dataset metadata

Resonance and Harmonics; Physical Model Grounding Reasoning, Multi-Formula Reasoning

\paragraph{Recorded answer/context.}
D: D: \$2 f\_0\$

\paragraph{Figure/image path.}
/root/proposal\_for\_physic/science-mango/phyx\_mini\_run/phyx\_data/test\_image/178.png

\paragraph{Formalization target.}
create a compiling Lean file with sorry bodies at `PhyXMiniProblems/problem\_phyx\_mini\_0178.lean`. The Lean declarations must preserve the physical quantities, dimensions or dimensional roles, figure labels, governing-law hypotheses, and final relation expressed by this problem.
Use Mathlib/Physlib names found through LeanExplore where available. If a physics API is missing, introduce faithful local abstractions rather than scalar placeholder aliases.

\begin{theorem}[Physics formalization target]
\label{thm:physics:phyx_mini_0178:target}
\lean{PhyXMiniProblems.ProblemPhyXMini0178.frequency_for_four_antinodes_after_fourfold_tension}
\uses{def:physics:phyx-mini-0178:phyxminiproblems-problemphyxmini0178-vibratingstringsetup, def:physics:phyx-mini-0178:phyxminiproblems-problemphyxmini0178-matchesinitialfigure, def:physics:phyx-mini-0178:phyxminiproblems-problemphyxmini0178-isrequiredfourantinodemode, def:physics:phyx-mini-0178:phyxminiproblems-problemphyxmini0178-tensionincreasedbyfour, def:physics:phyx-mini-0178:phyxminiproblems-problemphyxmini0178-obeystransversestringwavespeedlaw, def:physics:phyx-mini-0178:phyxminiproblems-problemphyxmini0178-obeysfixedendstandingwavelaw, def:physics:phyx-mini-0178:phyxminiproblems-problemphyxmini0178-haspositivephysicaldata}
The assigned autoformalize agent should translate this physics problem into Lean declarations in the covered file, with theorem and lemma proof bodies written as `by sorry`.
\end{theorem}
\begin{proof}
This is an autoformalization task, not a proof task. Produce statements that can later be proved by the physics prover without weakening the physical model.
\end{proof}

% --- Archon declaration topology begin ---
\section{Declaration topology}
\begin{definition}[String State]
  \label{def:physics:phyx-mini-0178:phyxminiproblems-problemphyxmini0178-stringstate}
  \lean{PhyXMiniProblems.ProblemPhyXMini0178.StringState}
  The two configurations compared in the problem
\end{definition}
\begin{proof}
  This is the declared model definition of String State.
\end{proof}
\begin{definition}[String Endpoint]
  \label{def:physics:phyx-mini-0178:phyxminiproblems-problemphyxmini0178-stringendpoint}
  \lean{PhyXMiniProblems.ProblemPhyXMini0178.StringEndpoint}
  The fixed supports visible at the left and right ends of the figure
\end{definition}
\begin{proof}
  This is the declared model definition of String Endpoint.
\end{proof}
\begin{definition}[Dim Length]
  \label{def:physics:phyx-mini-0178:phyxminiproblems-problemphyxmini0178-dimlength}
  \lean{PhyXMiniProblems.ProblemPhyXMini0178.DimLength}
  A physical length, independent of a choice of units
\end{definition}
\begin{proof}
  This is the declared model definition of Dim Length.
\end{proof}
\begin{definition}[Dim Frequency]
  \label{def:physics:phyx-mini-0178:phyxminiproblems-problemphyxmini0178-dimfrequency}
  \lean{PhyXMiniProblems.ProblemPhyXMini0178.DimFrequency}
  A physical (ordinary, not angular) frequency, with dimension inverse time
\end{definition}
\begin{proof}
  This is the declared model definition of Dim Frequency.
\end{proof}
\begin{definition}[Dim Tension]
  \label{def:physics:phyx-mini-0178:phyxminiproblems-problemphyxmini0178-dimtension}
  \lean{PhyXMiniProblems.ProblemPhyXMini0178.DimTension}
  String tension, carrying the physical dimension of force
\end{definition}
\begin{proof}
  This is the declared model definition of Dim Tension.
\end{proof}
\begin{definition}[Dim Linear Mass Density]
  \label{def:physics:phyx-mini-0178:phyxminiproblems-problemphyxmini0178-dimlinearmassdensity}
  \lean{PhyXMiniProblems.ProblemPhyXMini0178.DimLinearMassDensity}
  Linear mass density, carrying the physical dimension mass per length
\end{definition}
\begin{proof}
  This is the declared model definition of Dim Linear Mass Density.
\end{proof}
\begin{definition}[si Readout]
  \label{def:physics:phyx-mini-0178:phyxminiproblems-problemphyxmini0178-sireadout}
  \lean{PhyXMiniProblems.ProblemPhyXMini0178.siReadout}
  The real-valued SI readout of a nonnegative dimensionful quantity
\end{definition}
\begin{proof}
  This is the declared model definition of si Readout.
\end{proof}
\begin{definition}[Vibrating String Setup]
  \label{def:physics:phyx-mini-0178:phyxminiproblems-problemphyxmini0178-vibratingstringsetup}
  \lean{PhyXMiniProblems.ProblemPhyXMini0178.VibratingStringSetup}
  \uses{def:physics:phyx-mini-0178:phyxminiproblems-problemphyxmini0178-stringstate, def:physics:phyx-mini-0178:phyxminiproblems-problemphyxmini0178-stringendpoint, def:physics:phyx-mini-0178:phyxminiproblems-problemphyxmini0178-dimlength, def:physics:phyx-mini-0178:phyxminiproblems-problemphyxmini0178-dimfrequency, def:physics:phyx-mini-0178:phyxminiproblems-problemphyxmini0178-dimtension, def:physics:phyx-mini-0178:phyxminiproblems-problemphyxmini0178-dimlinearmassdensity}
  Data for the same string before and after its tension is increased. The boolean-free propositions record whether a configuration is a standing wave and whether each supported endpoint is a node
\end{definition}
\begin{proof}
  This is the declared model definition of Vibrating String Setup.
\end{proof}
\begin{definition}[Matches Initial Figure]
  \label{def:physics:phyx-mini-0178:phyxminiproblems-problemphyxmini0178-matchesinitialfigure}
  \lean{PhyXMiniProblems.ProblemPhyXMini0178.MatchesInitialFigure}
  \uses{def:physics:phyx-mini-0178:phyxminiproblems-problemphyxmini0178-vibratingstringsetup}
  The original diagram is a fixed-end standing wave with four antinodes
\end{definition}
\begin{proof}
  This is the declared model definition of Matches Initial Figure.
\end{proof}
\begin{definition}[Is Required Four Antinode Mode]
  \label{def:physics:phyx-mini-0178:phyxminiproblems-problemphyxmini0178-isrequiredfourantinodemode}
  \lean{PhyXMiniProblems.ProblemPhyXMini0178.IsRequiredFourAntinodeMode}
  \uses{def:physics:phyx-mini-0178:phyxminiproblems-problemphyxmini0178-vibratingstringsetup}
  The requested post-increase configuration is again a fixed-end standing wave with four antinodes
\end{definition}
\begin{proof}
  This is the declared model definition of Is Required Four Antinode Mode.
\end{proof}
\begin{definition}[Tension Increased By Four]
  \label{def:physics:phyx-mini-0178:phyxminiproblems-problemphyxmini0178-tensionincreasedbyfour}
  \lean{PhyXMiniProblems.ProblemPhyXMini0178.TensionIncreasedByFour}
  \uses{def:physics:phyx-mini-0178:phyxminiproblems-problemphyxmini0178-vibratingstringsetup}
  The stated fourfold increase of the physical string tension
\end{definition}
\begin{proof}
  This is the declared model definition of Tension Increased By Four.
\end{proof}
\begin{definition}[Obeys Transverse String Wave Speed Law]
  \label{def:physics:phyx-mini-0178:phyxminiproblems-problemphyxmini0178-obeystransversestringwavespeedlaw}
  \lean{PhyXMiniProblems.ProblemPhyXMini0178.ObeysTransverseStringWaveSpeedLaw}
  \uses{def:physics:phyx-mini-0178:phyxminiproblems-problemphyxmini0178-sireadout, def:physics:phyx-mini-0178:phyxminiproblems-problemphyxmini0178-vibratingstringsetup}
  The transverse-wave speed law for a stretched string, written without a square root as \textbackslash{}(v\textasciicircum{}2 μ = T\textbackslash{}) in SI readouts
\end{definition}
\begin{proof}
  This is the declared model definition of Obeys Transverse String Wave Speed Law.
\end{proof}
\begin{definition}[Obeys Fixed End Standing Wave Law]
  \label{def:physics:phyx-mini-0178:phyxminiproblems-problemphyxmini0178-obeysfixedendstandingwavelaw}
  \lean{PhyXMiniProblems.ProblemPhyXMini0178.ObeysFixedEndStandingWaveLaw}
  \uses{def:physics:phyx-mini-0178:phyxminiproblems-problemphyxmini0178-sireadout, def:physics:phyx-mini-0178:phyxminiproblems-problemphyxmini0178-vibratingstringsetup}
  For a standing wave with nodes at both fixed endpoints, \textbackslash{}(2 L f = n v\textbackslash{}), where \textbackslash{}(n\textbackslash{}) is the number of antinodes
\end{definition}
\begin{proof}
  This is the declared model definition of Obeys Fixed End Standing Wave Law.
\end{proof}
\begin{definition}[Has Positive Physical Data]
  \label{def:physics:phyx-mini-0178:phyxminiproblems-problemphyxmini0178-haspositivephysicaldata}
  \lean{PhyXMiniProblems.ProblemPhyXMini0178.HasPositivePhysicalData}
  \uses{def:physics:phyx-mini-0178:phyxminiproblems-problemphyxmini0178-sireadout, def:physics:phyx-mini-0178:phyxminiproblems-problemphyxmini0178-vibratingstringsetup}
  The nondegeneracy conditions for the physical string and both modes
\end{definition}
\begin{proof}
  This is the declared model definition of Has Positive Physical Data.
\end{proof}
\begin{definition}[Answer Choice]
  \label{def:physics:phyx-mini-0178:phyxminiproblems-problemphyxmini0178-answerchoice}
  \lean{PhyXMiniProblems.ProblemPhyXMini0178.AnswerChoice}
  Labels of the multiple-choice answers
\end{definition}
\begin{proof}
  This is the declared model definition of Answer Choice.
\end{proof}
\begin{definition}[answer Frequency Factor]
  \label{def:physics:phyx-mini-0178:phyxminiproblems-problemphyxmini0178-answerfrequencyfactor}
  \lean{PhyXMiniProblems.ProblemPhyXMini0178.answerFrequencyFactor}
  \uses{def:physics:phyx-mini-0178:phyxminiproblems-problemphyxmini0178-answerchoice}
  Dimensionless frequency multipliers displayed in choices A--D
\end{definition}
\begin{proof}
  This is the declared model definition of answer Frequency Factor.
\end{proof}
\begin{theorem}[wave Speed after fourfold tension]
  \label{thm:physics:phyx-mini-0178:phyxminiproblems-problemphyxmini0178-wavespeed-after-fourfold-tension}
  \lean{PhyXMiniProblems.ProblemPhyXMini0178.waveSpeed_after_fourfold_tension}
  \uses{def:physics:phyx-mini-0178:phyxminiproblems-problemphyxmini0178-vibratingstringsetup, def:physics:phyx-mini-0178:phyxminiproblems-problemphyxmini0178-tensionincreasedbyfour, def:physics:phyx-mini-0178:phyxminiproblems-problemphyxmini0178-obeystransversestringwavespeedlaw, def:physics:phyx-mini-0178:phyxminiproblems-problemphyxmini0178-haspositivephysicaldata}
  The wave-speed law and a fourfold tension increase make the wave speed double on the unchanged string
\end{theorem}
\begin{proof}
  The conclusion follows from the governing relations and figure readouts named in the statement of wave Speed after fourfold tension.
\end{proof}
% --- Archon declaration topology end ---
% --- Archon physics formalization source end ---
